\documentclass[11pt,leqno]{article}
\headheight=13.6pt

% packages
\usepackage[alphabetic]{amsrefs}
\usepackage{physics}
% margin spacing
\usepackage[top=1in, bottom=1in, left=0.5in, right=0.5in]{geometry}
\usepackage{amsfonts, amsmath, amssymb, amsthm}
\usepackage{extarrows}
\usepackage[none]{hyphenat}
\usepackage{fancyhdr}
\usepackage{graphicx}
\graphicspath{{./images/}}
\usepackage{float}
\usepackage{color}
\newcommand{\sai}[1]{\textcolor{red}{#1}}
\usepackage{enumitem}
% \usepackage{mathrsfs}
\usepackage{hyperref}
\usepackage[noabbrev, capitalise]{cleveref}
\crefformat{equation}{equation~#2#1#3}
\crefformat{lemma}{\textrm{Lemma}~#2#1#3}
\usepackage{quiver}

% theorems
\theoremstyle{plain}
\newtheorem{lem}{Lemma}
\newtheorem{lemma}[lem]{Lemma}
\newtheorem{thm}[lem]{Theorem}
\newtheorem{theorem}[lem]{Theorem}
\newtheorem{prop}[lem]{Proposition}
\newtheorem{proposition}[lem]{Proposition}
\newtheorem{cor}[lem]{Corollary}
\newtheorem{corollary}[lem]{Corollary}
\newtheorem{conj}[lem]{Conjecture}
\newtheorem{fact}[lem]{Fact}
\newtheorem{form}[lem]{Formula}

\theoremstyle{definition}
\newtheorem{defn}[lem]{Definition}
\newtheorem{definition/}[lem]{Definition}
\newenvironment{definition}
  {\renewcommand{\qedsymbol}{\textdagger}%
   \pushQED{\qed}\begin{definition/}}
  {\popQED\end{definition/}}
\newtheorem{example}[lem]{Example}
\newtheorem{remark}[lem]{Remark}
\newtheorem{exercise}[lem]{Exercise}
\newtheorem{notation}[lem]{Notation}

\numberwithin{equation}{section}
\numberwithin{lem}{section}

% header/footer formatting
\pagestyle{fancy}
\fancyhead{}
\fancyfoot{}
% \fancyhead[L]{Invariant theory and geometry}
\fancyhead[C]{}
% \fancyhead[R]{Sai Sivakumar}
% \fancyfoot[R]{\thepage}
\renewcommand{\headrulewidth}{0pt}

% paragraph indentation/spacing
\setlength{\parindent}{0cm}
\setlength{\parskip}{10pt}
\renewcommand{\baselinestretch}{1.25}

% mathscript S
\usepackage[mathscr]{euscript}

% math operators
\DeclareMathOperator{\Stab}{Stab}
\DeclareMathOperator{\Ind}{Ind}
\DeclareMathOperator{\Fun}{Fun}
\DeclareMathOperator{\id}{id}
\DeclareMathOperator{\End}{End}
\DeclareMathOperator{\GL}{GL}
\DeclareMathOperator{\SL}{SL}
\DeclareMathOperator{\SO}{SO}
\DeclareMathOperator{\Sp}{Sp}
\DeclareMathOperator{\Lie}{Lie}
\DeclareMathOperator{\diag}{diag}
\DeclareMathOperator{\lcm}{lcm}
\DeclareMathOperator{\U}{U}
\DeclareMathOperator{\SU}{SU}


% bracket notation for inner product
\usepackage{mathtools}

\DeclarePairedDelimiterX{\abr}[1]{\langle}{\rangle}{#1}

% smileys frownies
\usepackage{wasysym}
\newcommand{\smallhappy}{\raisebox{-.14em}{\smiley}}
\newcommand{\happy}{\raisebox{-.24em}{\resizebox{1.2em}{!}{\smiley}}}
\newcommand{\smallsad}{\raisebox{-.14em}{\frownie}}
\newcommand{\sad}{\raisebox{-.24em}{\resizebox{1.2em}{!}{\frownie}}}
\DeclareMathOperator{\mathhappy}{\!\happy\!}
\DeclareMathOperator{\smallmathhappy}{\!\smallhappy\!}
\DeclareMathOperator{\mathsad}{\!\sad\!}
\DeclareMathOperator{\smallmathsad}{\!\smallsad\!}

% set page count index to begin from 1
\setcounter{page}{1}

% array column and row separation
\arraycolsep = 2pt
\renewcommand{\arraystretch}{.6}

\begin{document}
\begin{enumerate}
    \item Projective GIT quotient
    \begin{enumerate}
        \item Projective varieties
        \begin{enumerate}
            \item $V = \mathbb C^{n+1}$, $\mathbb P(V) = (V\setminus\{0\})/\mathbb C^\times = \mathbb P^n$, $R = \mathbb C[V] = \mathbb C[x_0,\dots,x_n]$
            \item Homogeneous ideals of $R$ (i.e., ideals generated by homogeneous elements of $R$) define subsets of $\mathbb P^n$ called projective varieties
            \item A homogeneous radical ideal $I$ of $R$ gives rise to the subset $V(I) = \{[p] \mid f(p) = 0 \text{ for all } f\in I\}$ of $\mathbb P^n$; well defined since such $f$ are homogeneous polynomials
            \item the ``irrelevant ideal'' $R_+ = \bigoplus_{d>0}R_d = \abr{x_0,\dots,x_n}$ of positive degree homogeneous polynomials yields $V(R_+) = \emptyset$; we also have $V(\abr{a_ix_j - a_jx_i | 1\leq i,j\leq n}) = [a_0:\cdots:a_n]$
            \item Give $\mathbb P^n$ the Zariski topology where the closed sets are given by $V(I)$ for homogeneous ideals $I$
        \end{enumerate}
        \item Let a complex reductive group $G$ act linearly on $V$ (i.e., we are given a map $G \to \GL(V)$)
        \begin{enumerate}
            \item This action descends to an action on $\mathbb P^n$ since $G$ fixes zero and scalar multiplication commutes with the $G$-action
            \item The ring of invariants $R^G$ is graded and finitely generated by homogeneous polynomials; i.e., $R^G = \bigoplus_{d\geq 0}R^G_d$
            \item Define the open subset of semistable points (``semistable locus'') $\mathbb P^{ss}\coloneqq (\mathbb P^n)^{ss}$ by $\{[p] \mid f(p)\neq 0 \text{ for some } f\in R^G_d\}\subseteq P$; well defined since such $f$ are homogeneous polynomials
            \item The homogeneous generators $f_0,\dots,f_s$ of $R^G$ need not have the same degree; say $\deg(f_i) = d_i$
            \item With $d = \lcm(d_0,\dots,d_s)$; by some numerology (yoga of graded rings, see \href{https://stacks.math.columbia.edu/tag/0EGH}{this Stacks project lemma}), we should instead consider the generators $f_i^{D/d_i}$ of $R^G_{D}$, where $D=(s+1)d$ is the minimal multiple of the $d_i$ for which the following construction is guaranteed to work (we could do larger multiples if desired, or smaller multiples if the numerology works out in our favor)
            \item Let $I$ be the kernel of the map of graded rings $\mathbb C[t_0,\dots,t_s]\to \bigoplus_{\ell\geq 0}R^G_{\ell\cdot D}$ sending $t_i$ to $f_i^{D/d_i}$ (contrast with the affine GIT quotient!). The homogeneous ideal $I$ records the relations between the homogeneous generators $f_i^{D/d_i}$ of $R^G_{D}$, and this explains why it does not matter if we take $D$ to be a larger multiple of the $d_i$
        \end{enumerate}
        \item The projective GIT quotient of $\mathbb P^n$ by $G$ is given by the projection $\mathbb P^{ss}\to \mathbb P^n/\!/G\coloneqq V(I)\subseteq \mathbb P^s$ taking $[v]$ to $[f_0^{D/d_0}(v):\cdots:f_s^{D/d_s}(v)]$
        \begin{enumerate}
            \item We restrict to the subset of semistable points since we should not send $[v]$ to $[0\colon\cdots\colon 0]$
            \item The points of the variety $V(I)$ are in one to one correspondence with closures of $G$-orbits in the semistable locus $\mathbb P^{ss}$
            \item With $G = \mathbb C^\times$ acting on $\mathbb C^3$ by $t\cdot(x,y,z) = (tx,t^{-1}y,z)$, we obtain the action of $G$ on $\mathbb P^2$ by $t\cdot[x:y:z] = [tx:t^{-1}y:z]$, $R^G = \mathbb C[xy,z]$, $\mathbb P^{ss} = \mathbb P^2\setminus\{[1:0:0], [0:1:0]\}$, $D= 2\cdot \lcm(2,1) = 4$ $f_0^{D/d_0} = x^2y^2$, $f_1^{D/d_1} = z^4$, and $I = 0$ (there are no relations), so $\mathbb P^2/\!/\mathbb C^\times = V(I) = \mathbb P^1$
            \item Let $G= \mathbb C^\times$ act on $\mathbb C^4$ by $t\cdot(w,x,y,z) = (tw,tx,t^{-1}y,t^{-1}z)$, so that $R^G = \mathbb C[wy,wz,xy,xz]$. We would raise these generators to the fourth power ($D = 4\cdot \lcm(2,2,2,2) = 8$), but already we can see that the relations will be given by the defining ideal $I = \abr{t_0t_4-t_3t_2}$. Therefore $\mathbb P^3/\!/\mathbb C^\times = V(I) \cong \mathbb P^1\times \mathbb P^1\subseteq \mathbb P^3$, and we can see this since the map $[w:x:y:z]\mapsto [wy:wz:xy:xz]$ factors as $[w:x:y:z]\mapsto ([w:x],[y:z])\mapsto[wy:wz:xy:xz]$, where the second map is the Segre embedding
            \item In our original example with $\mathbb C^\times$ acting on $\mathbb C^2$ by $t\cdot (x,y) = (tx,ty)$, the affine quotient produces a copy of $\mathbb C$ whereas the projective quotient produces only a point. The qualitative difference is that in the copy of $\mathbb C$, each nonzero $a\in\mathbb C$ distinguishes each orbit $xy=a$ in $\mathbb C^2$ and $0\in\mathbb C$ corresponds to the orbits in $\mathbb C^2$ given by $x=0$, $y=0$, and $(0,0)$. The projective quotient forgets the latter orbits and identifies each hyperbola $xy=a$ to a point
            \item A variant of the first example where $\mathbb C^\times$ acts on $\mathbb C^3$ by $t\cdot(x,y,z) = (t^2x,y,tz)$ produces the same action on $\mathbb P^2$, but the invariants are different: $R^G$ is $\mathbb C[y]$, so the projective quotient is a point. So the data of the action of $G$ on $V$ is required
        \end{enumerate}
    \end{enumerate}
    \item Group actions on symplectic manifolds
    \begin{enumerate}
        \item Blackboxing symplectic theory
        \begin{enumerate}
            \item Given a linear action of $G$ on $V = \mathbb C^{n+1}$ ($n\geq 1$ to be safe), we may restrict the action to a maximal (analytic) compact subgroup $K$ of $G$
            \item For example, the usual maximal compacts we consider for the reductive groups 
            \[\GL_n(\mathbb C),\, \SL_n(\mathbb C),\, (\mathbb C^\times)^n,\, \SO_n(\mathbb C),\, \Sp_{2n}(\mathbb C)\]
            are
            \[\U(n),\, \SU(n),\, (S^1)^n,\, \SO_n(\mathbb R), \Sp_n(\mathbb H)\]
            \item A Hermitian inner product on $V$ may be used to define a symplectic $2$-form on $\mathbb P^n$ (as a real manifold). A symplectic form on an even-dimensional manifold is a smooth non-degenerate closed $2$-form
            \item By integrating over $K$, it is possible to choose a $K$-invariant Hermitian inner product on $V$ which then yields a $K$-invariant symplectic form on $\mathbb P^n$
            \item The only example we will care about in this talk is given by looking at what symplectic form the standard Hermitian inner product $\abr{-,-}$ on $V$ defines, the Fubini-Study form
        \end{enumerate}
        \item The Fubini-Study form $\omega_{\mathrm{FS}}$ on $\mathbb P^n$
        \begin{enumerate}
            \item We do something a little strange in order to express the Fubini-Study form using the usual Hermitian inner product on $V$
            \item For $[p]\in\mathbb P^n$, we can identify $T_{[p]}\mathbb P^n$ with the orthogonal space to $p$ in $\mathbb C^{n+1}$ with respect to the standard inner product on $\mathbb C^{n+1}$; this is via the map $d\pi_p\colon T_p(\mathbb C^{n+1}\setminus \{0\}) = \mathbb C^{n+1}\to T_{[p]}\mathbb P^n$ which takes $v$ to $\dv{t}[p+tv]|_{t=0}$; this map has kernel $\mathbb Cp$ which yields the earlier identification (maybe this is intuitive if we draw a picture and perturb a line in $V$). This orthogonal complement is an $n$-dimensional complex subspace of $V$
            \item Then the Fubini-Study form is given by the following formula: For $v,w\in T_{[p]}\mathbb P^n$, $\omega_{\mathrm{FS},[p]}(v,w) = \Im \abr{v,w}$ (we will take this as a definition)
            \item This description of the Fubini-Study form will be used to define something called a moment map which will satisfy certain properties
        \end{enumerate}
        \item The moment map corresponding to the action of the maximal compact $K$ subgroup of $G$ on $V$; we will for the sake of example stick to cases where $K$ is a closed subgroup of $\U(n+1)$, so that it preserves the usual inner product on $V$ (this is the case for all of the previous examples)
        \begin{enumerate}
            \item In general for a linear action of $K$ on $V$ preserving the inner product, the moment map we consider is the map $\mu\colon \mathbb P^n\to T_e^\ast K$ where for $A\in T_eK$, $\abr{\mu([p]), A}$ (the evaluation pairing) is $\abr{Ap,p}/(2i\norm{p}^2)$ (it's well defined, but maybe another way to think of this is that we can choose the representative $p$ to be on the unit sphere in $V$)
            \item We choose this map since it is compatible(?) with the symplectic structure in a specific way. Consider a tangent vector $v$ of $[p]$, which we think of as an element in the orthogonal complement of $p$ in $V$. Then the derivative of the map $\abr{\mu(-),A}$ at $[p]$ evaluated at $v$ is
            \[\dv{t}\abr{\mu(p+tv),A}|_{t=0} = \dv{t}\frac{\abr{A(p+tv),p+tv}}{2i\norm{p+tv}^2}\bigg|_{t=0} = \frac{1}{2i}(\abr{Ap,v}-\abr{v,Ap}) = \Im\abr{Ap,v},\]
            but this is just $\omega_{\mathrm{FS},[p]}(Ap,v)$. We should think of $Ap$ as the tangent vector of the infinitesimal flow along the action of $\exp At$ on $\mathbb P$, but technically, the vector $Ap$ should be the projection of $Ap$ onto the orthogonal complement of $p$ in $V$. There is no harm in using $Ap$ though since $\mathbb Cp$ is orthogonal to $v$
            (I would not carefully calculate this due to time constraints)
            \item The moment map is some kind of generalization of angular momentum; for example, the moment map on $T^\ast\mathbb R^3$ coming from lifting the natural action of $\SO_3$ on $\mathbb R^3$ is the cross product: pointwise it takes a position vector (the part in $\mathbb R^3$) and its momentum vector (the cotangent part) and calculates the cross product, which is the angular momentum about the origin
            \item For time, maybe I stick with one example; the case where $G = \mathbb C^\times$ acts on $\mathbb C^3$ by $t\cdot(x,y,z) = (tx,t^{-1}y,z)$ (so the action manifests $K = \mathbb S^1$ as a closed subgroup of $\U(n+1)$). We use the generator $A = i\diag(1,-1,0)$ of $T_eK$. Then the moment map we obtain is given explicitly by $\abr{\mu([p_0:p_1:p_2]), A} = \frac{\abr{i\diag(1,-1,0)(p_0,p_1,p_2), (p_0,p_1,p_2)}}{2i\norm{(p_0,p_1,p_2)}^2} = \frac{\abs{p_0}^2-\abs{p_1}^2}{2\norm{(p_0,p_1,p_2)}^2}$. 
        \end{enumerate}
    \end{enumerate}
    \item Kempf-Ness theorem
    \begin{enumerate}
        \item The statement is as follows: There is an inclusion $\mu^{-1}(0)\subseteq \mathbb P^{ss}$ which induces a homeomorphism of topological spaces $\mu^{-1}(0)/K\cong \mathbb P^n/\!/G$ (where the quotient on the left is an ordinary topological quotient)
        \item In our one example, $\mu^{-1}(0)$ is the locus $\{[p_0:p_1:p_2]\mid \abs{p_0} = \abs{p_1}\}$, and this really does include into the semistable locus $\mathbb P^{ss} = \mathbb P^2\setminus\{[1:0:0], [0:1:0]\}$
        \item Furthermore, the $S^1$ orbit of $\{[p_0:p_1:p_2]\mid \abs{p_0} = \abs{p_1}\}$ contains all points of the form $[sp_0:s^{-1}p_1:p_2]$, and these of course still belong to $\mu^{-1}(0)$
        \item The theorem says that the map taking $S^1[p_0:p_1:p_2]\in \mu^{-1}(0)/S^1$ to $[p_0^2p_1^2:p_2^4]\in\mathbb P^1 = \mathbb P^2/\!/\mathbb C^\times$ is a homeomorphism
        \item Should explain more the connection between the ``symplectic reduction'' appearing on the left hand side and the orbits parameterized on the right hand side of the Kempf-Ness homeomorphism
        \item What is the advantage of this theorem when it comes to studying (projective) GIT quotients?
        \item Why symplectic geometry/the moment map? It seems mysterious
    \end{enumerate}
\end{enumerate}

\end{document}